%!TEX root = ../example.tex
\chapter{Arquitectura del Sistema y Análisis de Tecnologías}
% Paquetes necesarios (deben estar en el preámbulo del documento principal)
% \usepackage{listings}
% \usepackage{amsmath}
% \usepackage{array}
% \usepackage{multirow}
% \usepackage{longtable}
% \usepackage{colortbl}
% \usepackage{xcolor}
% \usepackage{adjustbox}

\section{Introducción a la Arquitectura del Sistema}

El sistema de gestión de quejas y denuncias para la Defensoría del Pueblo (DDP) requiere una arquitectura robusta, segura y escalable que pueda manejar información sensible y garantizar la continuidad del servicio. Este capítulo presenta un análisis exhaustivo de las tecnologías seleccionadas, comparando alternativas y justificando cada decisión arquitectónica en función de los requisitos específicos del proyecto. Se abordarán no solo las características técnicas de cada componente, sino también su idoneidad para el contexto institucional, los requisitos normativos y la visión de largo plazo de la DDP.

\section{Marco Conceptual de la Arquitectura Propuesta}

La arquitectura propuesta sigue un modelo de microservicios distribuidos con un frontend en Angular y un backend híbrido que combina Java Spring Boot y Node.js. Esta combinación permite aprovechar las fortalezas de cada tecnología mientras se mantiene la cohesión y la integridad del sistema. La decisión de adoptar una arquitectura de microservicios no es trivial y responde a un análisis profundo de las necesidades actuales y futuras de la institución.

\subsection{Principios Arquitectónicos Fundamentales}

La arquitectura se fundamenta en los siguientes principios rectores:

\begin{itemize}
    \item \textbf{Seguridad por Diseño}: Cada capa de la arquitectura incorpora mecanismos de seguridad desde su concepción, no como una capa añadida posteriormente.
    \item \textbf{Resiliencia}: El sistema debe permanecer operativo incluso cuando componentes individuales fallen.
    \item \textbf{Escalabilidad Elástica}: Capacidad de crecer o reducir recursos según la demanda sin intervención manual.
    \item \textbf{Observabilidad}: Trazabilidad completa de las operaciones para auditoría y depuración.
    \item \textbf{Evolucionabilidad}: Capacidad de incorporar nuevas tecnologías sin reescribir el sistema completo.
    \item \textbf{Soberanía de Datos}: Los datos de los ciudadanos permanecen bajo control institucional.
\end{itemize}

\section{Análisis Comparativo de Frameworks Frontend}

El frontend es la cara visible del sistema para los funcionarios de la DDP y los ciudadanos. La selección del framework adecuado es crucial para garantizar una experiencia de usuario fluida, segura y accesible.

\begin{table}[htbp]
\centering
\caption{Análisis comparativo detallado de frameworks frontend}
\label{tab:comparativa-frontend}
\small
\begin{tabular}{|p{2.3cm}|p{3.3cm}|p{3.3cm}|p{3.3cm}|}
\hline
\textbf{Criterio} & \textbf{Angular} & \textbf{React} & \textbf{Vue.js} \\
\hline
\textbf{Seguridad} & \textbf{Alto}: TypeScript obligatorio, sanitización automática de contenido, protección XSS integrada, soporte nativo para CSP & \textbf{Medio}: JSX puede introducir vulnerabilidades, requiere librerías externas para sanitización & \textbf{Medio-Alto}: Plantillas con sanitización automática \\
\hline
\textbf{Arquitectura} & \textbf{Opinionada}: MVC completo, inyección de dependencias, módulos & \textbf{Basada en componentes}: Libertad total, requiere decisiones del equipo & \textbf{Flexible}: Progresiva, permite múltiples estilos \\
\hline
\textbf{Curva aprendizaje} & \textbf{Pronunciada}: TypeScript, RxJS, módulos & \textbf{Moderada}: JSX, hooks, estado & \textbf{Suave}: HTML templates, API simple \\
\hline
\textbf{Mantenibilidad} & \textbf{Excelente}: Estándares compartidos, estructura consistente & \textbf{Variable}: Depende de decisiones del equipo & \textbf{Buena}: Clara separación de concerns \\
\hline
\textbf{Escalabilidad organizacional} & \textbf{Alta}: Diseñado para equipos grandes, lazy loading & \textbf{Media}: Requiere configuración adicional & \textbf{Media-Buena}: Composition API mejora organización \\
\hline
\textbf{Soporte de pruebas} & \textbf{Nativo}: Karma, Jasmine, TestBed integrados & \textbf{Requiere configuración}: Jest/React Testing Library & \textbf{Requiere configuración}: Vue Test Utils \\
\hline
\end{tabular}
\end{table}

\subsection{Justificación Detallada de Angular}

Angular se selecciona como el framework frontend principal por las siguientes razones específicas para el contexto de la DDP:

\subsubsection{Seguridad Integrada y Cumplimiento Normativo}

La DDP maneja información sensible de ciudadanos, incluyendo datos personales, denuncias y documentación legal. Angular proporciona características de seguridad críticas:

\begin{itemize}
    \item \textbf{Sanitización automática}: Angular sanitiza automáticamente cualquier valor vinculado al DOM, previniendo ataques XSS (Cross-Site Scripting) sin configuración adicional. Esto es particularmente importante cuando se visualizan documentos y quejas que pueden contener contenido ingresado por usuarios malintencionados.
    
    \item \textbf{TypeScript obligatorio}: El tipado estricto reduce vulnerabilidades relacionadas con tipos dinámicos y facilita el análisis estático de código. Las vulnerabilidades de tipo son responsables del 15\% de las brechas de seguridad en aplicaciones JavaScript según estudios de la industria.
    
    \item \textbf{Soporte CSP}: Angular funciona nativamente con Content Security Policy, permitiendo políticas restrictivas sin romper funcionalidad. Esto permite a la DDP implementar las recomendaciones de OWASP para aplicaciones gubernamentales.
    
    \item \textbf{HTTP Interceptors}: Permiten implementar automáticamente tokens CSRF, headers de seguridad y manejo centralizado de errores, asegurando que todas las comunicaciones con el backend cumplan estándares de seguridad sin depender de la disciplina individual de cada desarrollador.
\end{itemize}

\subsubsection{Arquitectura Opinionada para Consistencia Institucional}

En una institución como la DDP, con múltiples equipos de desarrollo y posible rotación de personal, la arquitectura opinionada de Angular garantiza:

\begin{itemize}
    \item \textbf{Estandarización}: Todos los proyectos Angular siguen la misma estructura, facilitando la incorporación de nuevos desarrolladores y reduciendo el tiempo de onboarding.
    
    \item \textbf{Inyección de dependencias}: Sistema nativo que promueve el desacoplamiento y facilita las pruebas unitarias, permitiendo alcanzar coberturas de prueba superiores al 85\%.
    
    \item \textbf{Módulos funcionales}: Permite organizar el código por dominios de negocio (quejas, usuarios, reportes) con lazy loading para optimizar rendimiento. Esto significa que los funcionarios solo cargan el código necesario para su función específica.
    
    \item \textbf{CLI robusto}: Angular CLI genera código consistente, aplica mejores prácticas y automatiza tareas de build y testing, reduciendo errores humanos.
\end{itemize}

\section{Análisis Comparativo de Tecnologías Backend}

El backend requiere una combinación de tecnologías que ofrezcan robustez para operaciones críticas y flexibilidad para capas de integración.

\begin{table}[htbp]
\centering
\caption{Análisis comparativo detallado de tecnologías backend}
\label{tab:comparativa-backend}
\small
\begin{tabular}{|p{2.5cm}|p{4cm}|p{4cm}|}
\hline
\textbf{Criterio} & \textbf{Java Spring Boot} & \textbf{Node.js} \\
\hline
\textbf{Seguridad} & \textbf{Muy alto}: Spring Security, autenticación JWT/OAuth2, autorización basada en roles, protección CSRF, auditoría integrada & \textbf{Alto}: Helmet.js, express-rate-limit, bcrypt, JWT, requiere configuración explícita \\
\hline
\textbf{Rendimiento} & \textbf{Excelente para operaciones complejas}: Multi-threading, ideal para operaciones transaccionales y computación intensiva & \textbf{Excelente para I/O}: Modelo asíncrono no bloqueante, ideal para APIs REST simples y tiempo real \\
\hline
\textbf{Escalabilidad} & \textbf{Alta}: Soporte nativo para microservicios con Spring Cloud, service discovery & \textbf{Alta}: Cluster module, PM2, arquitectura orientada a eventos \\
\hline
\textbf{Madurez} & \textbf{Muy alta}: 20+ años en entornos empresariales, utilizado por gobiernos y bancos & \textbf{Alta}: 10+ años, ampliamente adoptado \\
\hline
\textbf{Transaccionalidad} & \textbf{Soporte completo}: Transactional, gestión declarativa de transacciones & \textbf{Limitado}: No tiene transacciones nativas, requiere manejo manual \\
\hline
\end{tabular}
\end{table}

\subsection{Análisis Profundo de Java y Spring Boot}

\subsubsection{Fundamentos de Java en el Contexto Gubernamental}

Java es uno de los lenguajes más utilizados en el sector público a nivel mundial por razones fundamentales:

\begin{itemize}
    \item \textbf{Madurez probada}: Presente en sistemas críticos desde hace más de 25 años, incluyendo aplicaciones bancarias, gubernamentales y de defensa.
    
    \item \textbf{Estandarización}: Procesos JCP (Java Community Process) que garantizan evolución controlada y retrocompatibilidad.
    
    \item \textbf{Rendimiento predecible}: La JVM (Java Virtual Machine) ofrece un rendimiento consistente bajo carga, con garbage collectors configurables para diferentes perfiles.
    
    \item \textbf{Ecosistema de herramientas}: Amplio soporte de herramientas de monitoreo, profiling y debugging como JProfiler, VisualVM, y Java Flight Recorder.
\end{itemize}

\subsubsection{Spring Boot como Marco de Desarrollo Empresarial}

Spring Boot extiende las capacidades de Java con características específicas para aplicaciones modernas:

\textbf{Inyección de Dependencias y IoC (Inversión de Control)}:

\begin{lstlisting}[language=Java, caption=Ejemplo de inyección de dependencias en Spring Boot]
@Service
public class QuejaService {
    
    private final QuejaRepository quejaRepository;
    private final NotificacionService notificacionService;
    
    public QuejaService(QuejaRepository quejaRepository, 
                        NotificacionService notificacionService) {
        this.quejaRepository = quejaRepository;
        this.notificacionService = notificacionService;
    }
    
    @Transactional
    public Queja registrarQueja(QuejaDTO quejaDTO) {
        Queja queja = mapper.toEntity(quejaDTO);
        queja = quejaRepository.save(queja);
        notificacionService.notificarRegistro(queja);
        return queja;
    }
}
\end{lstlisting}

\textbf{Spring Security y Control de Acceso}:

\begin{lstlisting}[language=Java, caption=Configuración de seguridad para la DDP]
@Configuration
@EnableWebSecurity
public class SecurityConfig {
    
    @Bean
    public SecurityFilterChain filterChain(HttpSecurity http) throws Exception {
        http
            .authorizeRequests()
                .antMatchers("/api/quejas/**").hasAnyRole("ABOGADO", "DEFENSOR", "ADMIN")
                .antMatchers("/api/admin/**").hasRole("ADMIN")
                .antMatchers("/api/reportes/**").hasAnyRole("ANALISTA", "ADMIN")
                .anyRequest().authenticated()
            .and()
            .oauth2ResourceServer()
                .jwt()
            .and()
            .sessionManagement()
                .sessionCreationPolicy(SessionCreationPolicy.STATELESS);
        
        return http.build();
    }
}
\end{lstlisting}

\subsection{Node.js como Complemento Estratégico}

Node.js complementa la arquitectura en roles específicos donde su modelo asíncrono ofrece ventajas:

\subsubsection{API Gateway}

\begin{itemize}
    \item \textbf{Alto rendimiento I/O}: Maneja miles de conexiones concurrentes con bajo consumo de recursos, ideal para el punto de entrada del sistema.
    
    \item \textbf{Proxy inverso}: Express/Fastify como proxy para enrutar peticiones a microservicios, reduciendo latencia.
    
    \item \textbf{Rate limiting}: Implementación eficiente de límites de peticiones para prevenir ataques de denegación de servicio.
\end{itemize}

\textbf{Estructura conceptual de API Gateway}:

\begin{lstlisting}[caption=Estructura conceptual de API Gateway]
// El API Gateway actua como punto unico de entrada
// Redirige peticiones a los microservicios correspondientes
// Implementa rate limiting, autenticacion y logging centralizado

Funcion principal del Gateway:
- Recibir peticiones del cliente
- Validar tokens JWT
- Aplicar limites de tasa
- Enrutar al servicio correspondiente
- Agregar respuestas cuando sea necesario
- Registrar todas las operaciones para auditoria
\end{lstlisting}

\subsubsection{Servicio de Notificaciones}

\begin{itemize}
    \item \textbf{Tiempo real}: Socket.io para notificaciones en tiempo real a funcionarios cuando se asignan nuevas quejas.
    
    \item \textbf{Email/SMS}: Nodemailer, Twilio, librerías maduras para comunicación multicanal con bajo overhead.
    
    \item \textbf{WebSockets}: Manejo eficiente de conexiones persistentes para actualizaciones de estado.
\end{itemize}

\section{Análisis del Sistema de Base de Datos: PostgreSQL}

La selección del motor de base de datos es crítica para garantizar integridad, seguridad y rendimiento.

\begin{table}[htbp]
\centering
\caption{Análisis exhaustivo de PostgreSQL para el proyecto DDP}
\label{tab:postgresql-analisis}
\small
\begin{tabular}{|p{3cm}|p{4.5cm}|p{3.8cm}|}
\hline
\textbf{Aspecto} & \textbf{Ventajas} & \textbf{Consideraciones} \\
\hline
\textbf{Seguridad} & Cumplimiento ACID, cifrado SSL/TLS, roles y permisos granulares, Row-Level Security, auditoría con pgAudit & Requiere configuración adecuada de permisos \\
\hline
\textbf{Integridad} & Foreign keys, constraints, transacciones ACID, triggers, vistas materializadas & Modelo relacional estricto \\
\hline
\textbf{Madurez} & 30+ años de desarrollo, comunidad global, documentación extensa & - \\
\hline
\textbf{Rendimiento} & Optimizador avanzado, índices múltiples, particionamiento nativo, replicación streaming & Requiere tuning periódico \\
\hline
\textbf{JSONB soporte} & Datos semiestructurados, índices GIN sobre JSONB & Complementa sin perder relacional \\
\hline
\end{tabular}
\end{table}

\subsection{Justificación Detallada de PostgreSQL}

\subsubsection{Integridad de Datos Jurídicos}

La información manejada por la DDP (quejas, denuncias, expedientes) requiere garantías absolutas de integridad:

\begin{itemize}
    \item \textbf{Transacciones atómicas}: Las operaciones que involucran múltiples tablas se ejecutan completamente o no se ejecutan, garantizando que no existan estados inconsistentes.
    
    \item \textbf{Integridad referencial}: Las foreign keys garantizan que no existan huérfanos. Si una queja hace referencia a un ciudadano, ese ciudadano debe existir.
    
    \item \textbf{Check constraints}: Validaciones de negocio a nivel base de datos, como garantizar que la fecha de resolución sea posterior a la fecha de registro.
\end{itemize}

\begin{lstlisting}[language=SQL, caption=Constraints de integridad en PostgreSQL]
CREATE TABLE quejas (
    id SERIAL PRIMARY KEY,
    ciudadano_id INTEGER NOT NULL REFERENCES ciudadanos(id),
    fecha_registro TIMESTAMP NOT NULL DEFAULT CURRENT_TIMESTAMP,
    fecha_resolucion TIMESTAMP,
    estado VARCHAR(20) NOT NULL CHECK (estado IN ('REGISTRADA', 'EN_TRAMITE', 'RESUELTA')),
    CONSTRAINT fecha_resolucion_valida CHECK (
        fecha_resolucion IS NULL OR fecha_resolucion >= fecha_registro
    )
);
\end{lstlisting}

\section{Análisis Comparativo de Arquitecturas: Microservicios vs Monolito}

La decisión arquitectónica fundamental define cómo evolucionará el sistema en los próximos años.

\begin{table}[htbp]
\centering
\caption{Comparativa detallada: Microservicios vs Arquitectura Monolítica}
\label{tab:comparativa-arquitectura}
\small
\begin{tabular}{|p{2.8cm}|p{4.2cm}|p{4.2cm}|}
\hline
\textbf{Criterio} & \textbf{Microservicios} & \textbf{Monolito} \\
\hline
\textbf{Escalabilidad} & Independiente por servicio, escalado granular, optimización por carga & Escalado completo, desperdicio de recursos, cuello de botella único \\
\hline
\textbf{Mantenibilidad} & Código modular por dominio, despliegue independiente, equipos autónomos & Código acoplado, cambios afectan toda la aplicación, coordinación requerida \\
\hline
\textbf{Seguridad} & Aislable por servicio, menor superficie de ataque, segmentación de red & Centralizada, mayor superficie de ataque, una brecha expone todo \\
\hline
\textbf{Tolerancia a fallos} & Circuit breakers, aislamiento de fallos, degradación graceful & Fallo total, sin aislamiento, caída completa del sistema \\
\hline
\textbf{Complejidad operacional} & Mayor: múltiples servicios, requiere orquestación, trazabilidad distribuida & Menor: un solo despliegue, monitoreo simple, logs centralizados \\
\hline
\end{tabular}
\end{table}

\subsection{Análisis desde la Perspectiva Analítica}

Desde un punto de vista analítico, la decisión de adoptar microservicios se fundamenta en la separación por dominios funcionales:

\begin{enumerate}
    \item \textbf{Servicio de Gestión de Quejas}: Core del sistema, maneja el ciclo de vida completo de quejas y denuncias. Requiere alta consistencia transaccional.
    
    \item \textbf{Servicio de Autenticación y Autorización}: Gestiona usuarios, roles y permisos. Debe ser extremadamente seguro y auditado.
    
    \item \textbf{Servicio de Documentos}: Almacena y gestiona archivos adjuntos. Maneja grandes volúmenes de datos binarios.
    
    \item \textbf{Servicio de Notificaciones}: Envía alertas por email, SMS, notificaciones push. Es intensivo en I/O.
    
    \item \textbf{Clasificador IA}: Procesa quejas entrantes para clasificarlas automáticamente. Requiere recursos computacionales específicos.
    
    \item \textbf{Servicio de Reportes}: Genera estadísticas y reportes para la gestión.
\end{enumerate}

\subsection{Análisis desde la Perspectiva Técnica}

\subsubsection{Patrones de Resiliencia Implementados}

\begin{lstlisting}[language=Java, caption=Ejemplo de Circuit Breaker con Resilience4j]
@Service
public class ClasificadorService {
    
    private final ClasificadorClient clasificadorClient;
    
    @CircuitBreaker(name = "clasificadorIA", 
                     fallbackMethod = "clasificarFallback")
    public Clasificacion clasificarQueja(Queja queja) {
        return clasificadorClient.clasificar(queja);
    }
    
    public Clasificacion clasificarFallback(Queja queja, Exception ex) {
        log.warn("Clasificador IA no disponible, usando clasificacion por defecto", ex);
        return new Clasificacion("PENDIENTE", "GENERAL");
    }
}
\end{lstlisting}

\subsection{Análisis desde la Perspectiva de Visión a Futuro}

\subsubsection{Evolución Independiente de Componentes}

La arquitectura de microservicios permite que el sistema evolucione orgánicamente:

\begin{itemize}
    \item \textbf{Actualizaciones tecnológicas}: Se puede migrar un servicio de Spring Boot a otra tecnología sin reescribir todo el sistema.
    
    \item \textbf{Incorporación de nuevas funcionalidades}: Nuevos dominios de negocio pueden agregarse como servicios independientes.
    
    \item \textbf{Adaptación a cambios normativos}: Actualizar rápidamente un servicio para cumplir nuevas regulaciones sin afectar otros.
\end{itemize}

\subsection{Desafíos y Mitigaciones de Microservicios}

\subsubsection{Complejidad de Transacciones Distribuidas}

Las transacciones que abarcan múltiples servicios requieren patrones específicos:

\begin{itemize}
    \item \textbf{SAGA Coreográfica}: Cada servicio ejecuta su transacción y publica eventos.
    \item \textbf{SAGA Orquestada}: Un orquestador central coordina la transacción distribuida.
    \item \textbf{Eventual consistency}: Se acepta consistencia eventual para operaciones no críticas.
\end{itemize}

\begin{lstlisting}[language=Java, caption=Patrón SAGA orquestado]
@Component
public class RegistroQuejaOrquestador {
    
    @Autowired
    private QuejaService quejaService;
    
    @Autowired
    private DocumentoService documentoService;
    
    public void registrarQuejaCompleta(QuejaDTO quejaDTO) {
        try {
            // Paso 1: Registrar queja
            Queja queja = quejaService.registrar(quejaDTO);
            
            // Paso 2: Adjuntar documentos
            documentoService.adjuntar(queja.getId(), quejaDTO.getDocumentos());
            
        } catch (Exception e) {
            // Compensar transacciones
            compensar(e, quejaDTO);
        }
    }
    
    private void compensar(Exception e, QuejaDTO quejaDTO) {
        log.error("Error en registro de queja, compensando...", e);
        // Logica de compensacion
    }
}
\end{lstlisting}

\section{Matriz de Selección Tecnológica y Evaluación de Riesgos}

\begin{table}[htbp]
\centering
\caption{Matriz de selección tecnológica con ponderación por criterios}
\label{tab:matriz-seleccion}
\small
\begin{tabular}{|p{2.2cm}|c|c|c|c|c|c|}
\hline
\textbf{Tecnología} & \textbf{Seguridad (30\%)} & \textbf{Escalabilidad (20\%)} & \textbf{Mantenibilidad (20\%)} & \textbf{Curva (10\%)} & \textbf{Estabilidad (20\%)} & \textbf{Puntuación} \\
\hline
Angular & 5 (1.5) & 4 (0.8) & 5 (1.0) & 3 (0.3) & 5 (1.0) & \textbf{4.6} \\
\hline
React & 3 (0.9) & 4 (0.8) & 3 (0.6) & 4 (0.4) & 4 (0.8) & \textbf{3.5} \\
\hline
Vue.js & 4 (1.2) & 4 (0.8) & 4 (0.8) & 5 (0.5) & 4 (0.8) & \textbf{4.1} \\
\hline
Spring Boot & 5 (1.5) & 5 (1.0) & 5 (1.0) & 3 (0.3) & 5 (1.0) & \textbf{4.8} \\
\hline
Node.js & 4 (1.2) & 5 (1.0) & 4 (0.8) & 4 (0.4) & 4 (0.8) & \textbf{4.2} \\
\hline
PostgreSQL & 5 (1.5) & 4 (0.8) & 4 (0.8) & 3 (0.3) & 5 (1.0) & \textbf{4.4} \\
\hline
Microservicios & 4 (1.2) & 5 (1.0) & 5 (1.0) & 2 (0.2) & 4 (0.8) & \textbf{4.2} \\
\hline
Monolito & 4 (1.2) & 3 (0.6) & 3 (0.6) & 4 (0.4) & 4 (0.8) & \textbf{3.6} \\
\hline
\end{tabular}
\end{table}

\subsection{Análisis de Riesgos Tecnológicos}

\begin{table}[htbp]
\centering
\caption{Matriz de riesgos tecnológicos y mitigaciones}
\label{tab:riesgos}
\small
\begin{tabular}{|p{3.2cm}|p{2.5cm}|p{5cm}|}
\hline
\textbf{Riesgo} & \textbf{Impacto} & \textbf{Mitigación} \\
\hline
\textbf{Complejidad de microservicios} & Alto & Service mesh, OpenTelemetry para trazabilidad, Kubernetes para orquestación \\
\hline
\textbf{Seguridad en Node.js} & Medio & SCA (Snyk/npm audit), actualización automática, análisis estático \\
\hline
\textbf{Rendimiento de PostgreSQL} & Medio & Monitoreo, índices adecuados, particionamiento, pooling con PgBouncer \\
\hline
\textbf{Obsolescencia tecnológica} & Bajo & Versiones LTS, contenerización, actualizaciones programadas \\
\hline
\textbf{Brecha de conocimiento} & Medio & Programa de capacitación, consultoría inicial, documentación \\
\hline
\end{tabular}
\end{table}

\section{Estrategia de Implementación y Roadmap Tecnológico}

\subsection{Fase 1: Fundamentos (Meses 1-3)}
\begin{itemize}
    \item Implementación del API Gateway con Node.js
    \item Servicio de Autenticación y Autorización con Spring Boot
    \item Base de datos PostgreSQL
    \item Despliegue en contenedores Docker
\end{itemize}

\subsection{Fase 2: Core del Sistema (Meses 4-6)}
\begin{itemize}
    \item Servicio de Gestión de Quejas con Spring Boot
    \item Frontend básico en Angular para funcionarios
    \item Integración de RabbitMQ para comunicación asíncrona
\end{itemize}

\subsection{Fase 3: Funcionalidades Avanzadas (Meses 7-9)}
\begin{itemize}
    \item Servicio de Documentos
    \item Clasificador IA con Python/Flask
    \item Servicio de Notificaciones con Node.js
\end{itemize}

\subsection{Fase 4: Optimización (Meses 10-12)}
\begin{itemize}
    \item Kubernetes para orquestación avanzada
    \item Caché distribuida con Redis
    \item Particionamiento de bases de datos
    \item Alta disponibilidad
\end{itemize}

\section{Conclusiones de la Selección Tecnológica}

La combinación tecnológica seleccionada para el sistema de la DDP responde a las necesidades específicas de una institución pública que maneja información sensible y requiere:

\begin{itemize}
    \item \textbf{Seguridad}: Angular y Spring Boot proporcionan características de seguridad maduras y probadas en entornos gubernamentales.
    
    \item \textbf{Escalabilidad}: Microservicios permiten crecimiento independiente por dominio funcional.
    
    \item \textbf{Mantenibilidad}: La arquitectura modular garantiza consistencia a largo plazo.
    
    \item \textbf{Integridad de datos}: PostgreSQL asegura cumplimiento ACID y trazabilidad legal.
    
    \item \textbf{Flexibilidad}: Node.js complementa para casos donde su modelo asíncrono es ventajoso.
    
    \item \textbf{Visión a futuro}: La arquitectura prepara a la DDP para un crecimiento sostenible, adaptándose a futuros requisitos sin comprometer la seguridad ni la integridad de los datos.
\end{itemize}

Esta arquitectura establece una base sólida para la transformación digital de la institución, garantizando que pueda cumplir su misión con herramientas tecnológicas modernas, seguras y confiables.